% VI/48 reader bridge -- additive exposition, original exercises unchanged.
\paragraph{Why a short exact sequence is not immediately exact on sections.}
Let $X$ be a scheme and
\[
  0\longrightarrow\mathcal F'\xrightarrow{\alpha}\mathcal F
   \xrightarrow{\beta}\mathcal F''\longrightarrow0
\]
a short exact sequence of sheaves of abelian groups (or modules).
The surjection $\beta$ means that sections of $\mathcal F''$ admit
\emph{local} lifts after possibly shrinking an open set.  It need
not mean that every section over all of $X$ has a global lift.
Taking global sections always yields a left-exact sequence
\[
0\longrightarrow\Gamma(X,\mathcal F')
 \longrightarrow\Gamma(X,\mathcal F)
 \longrightarrow\Gamma(X,\mathcal F''),
\]
but surjectivity on the right is a separate question.  Its first
obstruction is measured by $H^1(X,\mathcal F')$.

\paragraph{The connecting map, constructed from local lifts.}
Take an open cover $\mathfrak U=(U_i)$ and a global section
$s\in\Gamma(X,\mathcal F'')$.  Refine the cover, if necessary,
until there are lifts $t_i\in\mathcal F(U_i)$ with
$\beta(t_i)=s|_{U_i}$.  On $U_i\cap U_j$ the difference
$t_j-t_i$ maps to zero under $\beta$, so it is the image of a
unique section
\[
 c_{ij}\in\mathcal F'(U_i\cap U_j),
 \qquad\alpha(c_{ij})=t_j-t_i.
\]
The relation
$(t_j-t_i)+(t_k-t_j)+(t_i-t_k)=0$ on triple intersections
makes $(c_{ij})$ a \v{C}ech $1$-cocycle.  Altering the lifts
$t_i$ changes it by a coboundary.  On an acyclic cover, or
after passing to a suitable refinement in the general
construction, this gives the connecting homomorphism
\[
 \delta:H^0(X,\mathcal F'')\longrightarrow H^1(X,\mathcal F').
\]
The class $\delta(s)$ vanishes precisely when $s$ is the image
of a global section of $\mathcal F$.  This is an explicit
version of the obstruction represented by the long exact
cohomology sequence.

\paragraph{The long exact sequence and its exactness test.}
The derived-functor sequence has the form
\[
\begin{aligned}
0&\to H^0(\mathcal F')\to H^0(\mathcal F)\to H^0(\mathcal F'')
 \xrightarrow{\delta} H^1(\mathcal F')\\
 &\to H^1(\mathcal F)\to H^1(\mathcal F'')
 \xrightarrow{\delta}H^2(\mathcal F')\to\cdots .
\end{aligned}
\]
(All cohomology groups in this display are taken on $X$.)
At every group, the image of the incoming arrow equals the
kernel of the outgoing arrow.  In particular, if
$H^1(X,\mathcal F')=0$, then
$\Gamma(X,\mathcal F)\to\Gamma(X,\mathcal F'')$ is surjective.
Conversely, surjectivity of that one map does \emph{not}
force $H^1(X,\mathcal F')$ to vanish in general.

\paragraph{Worked example 1: the Euler sequence on $\mathbb P^1$.}
Let $X=\mathbb P^1_k$ with homogeneous coordinates $[X_0:X_1]$.
The maps
\[
\begin{aligned}
0\longrightarrow\mathcal O(-2)
 &\xrightarrow{\ \alpha\ }\mathcal O(-1)^{\oplus2}
 \xrightarrow{\ \beta\ }\mathcal O\longrightarrow0,\\
\alpha(g)&=(-X_1g,X_0g),\\
\beta(a,b)&=X_0a+X_1b
\end{aligned}
\]
give a short exact sequence.  To check it, cover $X$ by
$U_0=D_+(X_0)$ and $U_1=D_+(X_1)$.  Where $X_0$ is invertible,
the section $1$ of $\mathcal O$ lifts to $(X_0^{-1},0)$;
where $X_1$ is invertible it lifts to $(0,X_1^{-1})$.
The formulas also show that the kernel is generated locally
by $(-X_1,X_0)$, giving the displayed identification with
$\mathcal O(-2)$.

The two lifts of $1$ do not agree on the overlap.  Their
difference, second minus first, is
\[
(-X_0^{-1},X_1^{-1})
 =\alpha\!\left(\frac{1}{X_0X_1}\right).
\]
Write $t=X_1/X_0$.  Relative to the standard frame
$X_0^{-2}$ of $\mathcal O(-2)$ on $U_0$, the cocycle
$(X_0X_1)^{-1}$ has coefficient $t^{-1}$.  Hence
\[
 \delta(1)=[t^{-1}]
 \in H^1(\mathbb P^1_k,\mathcal O(-2))\simeq k.
\]
The calculation of VI/47 shows that this class is nonzero:
\[
H^0(\mathcal O(-1))=0,\qquad
H^1(\mathcal O(-1))=0,\qquad
H^1(\mathcal O(-2))=k.
\]
Thus $\delta:k\to k$ is an isomorphism.  A sheaf surjection
onto $\mathcal O$ need not be surjective on global sections.

\paragraph{Worked example 2: a point-evaluation exact sequence.}
Fix a $k$-rational point $P\in\mathbb P^1_k$, and denote its
skyscraper sheaf by $k(P)$.  Its ideal sheaf is
$\mathcal I_P\simeq\mathcal O(-1)$, giving
\[
0\to\mathcal O(-1)\to\mathcal O\to k(P)\to0.
\]
On global sections this becomes
$0\to0\to k\xrightarrow{\sim}k\to0$, since the constant
section $a$ restricts to $a$ at $P$.  There is no nonzero
connecting class here: $H^1(\mathcal O(-1))=0$.
This contrasts with the preceding Euler sequence and shows
why the source sheaf, not merely the existence of a short
exact sequence, determines whether lifting is obstructed.

\paragraph{Warning about passing to \v{C}ech complexes.}
A short exact sequence of sheaves need not give a
term-by-term short exact sequence after applying
$\mathcal F\mapsto\mathcal F(U_{i_0}\cap\cdots\cap U_{i_p})$:
restriction to a fixed open set is left exact on sections,
not generally right exact.  For the quasi-coherent sheaves
in the examples above, the chosen affine intersections
are acyclic and taking sections there is exact.  In the
general case, derived functors (or suitable resolutions)
give the long exact sequence; one must not invoke the
snake lemma on arbitrary \v{C}ech cochain groups without
checking this hypothesis.
